% problem-000148 7 有机酸性与反应性
\section*{7. 有机酸性与反应性}
\textit{下列为分问。}

\subsection*{7-1}
酸碱理论对有机化学的发展具有重要的意义。

\subsection*{7-1}
-1 苯甲酸衍⽣物的酸性与其取代基的性质紧密相关。如下所⽰，化合物B的 ${ \mathfrak { o } } K _ { a }$ 值显著⼩于其他三个化合物。对此结果给出合理解释。

（图：）

\subsection*{7-1}
-2 ⼆元羧酸的酸性与⼆个羧基的距离密切相关。例如，草酸的 $\mathsf { p } K _ { a 1 }$ ⼩于丙⼆酸的 $\mathsf { p } K _ { a 1 }$ ，邻苯⼆甲酸的 $\mathcal { K } _ { a \mathrm { : } }$ /⼩于间苯⼆甲酸的 $\mathsf { p } K _ { a 1 } \mathfrak { c }$ 。对此结果给出合理解释。

\subsection*{7-1}
-3 2,6-⼆甲基-4-硝基苯酚的 $\mathsf { p } K _ { a }$ ⼩于 3,5-⼆甲基-4-硝基苯酚的 $\mathsf { p } K _ { a }$ ，对此结果给出合理解释。

\subsection*{7-2}
2015年，英国的科研团队⾸次合成了全顺式的1,2,3,4,5,6-六氟环⼰烷(1)。这是迄今已知极性最⼤的⾮离⼦型有机化合物。1的合成路线如下图所⽰：以肌醇(myo-inositol, 2)为起始原料，通过6步转化合成中间体3，再经过以脱氧氟化和亲核氟化为关键步骤的6步反应制备1。(提⽰：DeoxoFluor是⼀种常⽤的脱氧氟化试剂，可将羟基转化为氟，并⽣成 $\mathrm { S O F _ { 2 c } }$ 。Et_{3}N·3HF是⼀种常⽤的亲核氟化试剂)

（图：）

\subsection*{7-2}
-1 1,2,3,4,5,6-六氟环已烷有多少个⽴体异构体（仅需数⽬）？画出其中⼀对对映体的结构简式。

\subsection*{7-2}
-2 画出化合物4的结构简式。

\subsection*{7-2}
-3 检测表明化合物 1 的 F-NMR 谱图只有−212.8 ppm 和−217.3 ppm 两种信号。画出 1 的稳定构象，并解释它具有⼤极性的原因。

\subsection*{7-2}
-4 简要阐述最后⼀步反应（由8合成1）产率低的主要原因。
